% Pandoc template for IEEE conference paper sections
\section{Improved Cocks-Pinch with NTT Constraints}

We modify the traditional Cocks-Pinch algorithm to guarantee high
two-adicity in both output primes. The key insight is that the algebraic
structure of the cyclotomic polynomial \(\Phi_{18}\) allows
deterministic control over the two-adicity of \(r\), while an extended
lift search achieves the same for \(p\).

\subsection{NTT-Friendly Scalar Field via Seed
Constraint}\label{ntt-friendly-scalar-field-via-seed-constraint}

Since \(r = \Phi_{18}(T) = T^6 - T^3 + 1\), we observe:

\[r - 1 = T^6 - T^3 = T^3(T^3 - 1)\]

If \(T \equiv 0 \pmod{2^m}\), then \(T^3 \equiv 0 \pmod{2^{3m}}\). Since
\(T^3 - 1\) is odd (as \(T^3\) is even), we obtain:

\[\nu_2(r - 1) = \nu_2(T^3) = 3m\]

\textbf{Consequence:} To guarantee \(\nu_2(r-1) \geq s\), it suffices to
choose \(T \equiv 0 \pmod{2^{\lceil s/3 \rceil}}\). This is a
\emph{deterministic} guarantee requiring no rejection sampling.

For our target of \(s = 36\), we set \(m = \lceil 36/3 \rceil = 12\),
constraining \(T \equiv 0 \pmod{2^{12}}\). The search space remains
\(2^{86-12} = 2^{74}\) candidates---more than sufficient for finding
prime \(r\).

\subsection{NTT-Friendly Base Field via Extended Lift
Search}\label{ntt-friendly-base-field-via-extended-lift-search}

After obtaining a valid \(r\) with the desired two-adicity, the base
field prime is:

\[p = \frac{(t_0 + h_t r)^2 + 3(y_0 + h_y r)^2}{4}\]

The two-adicity of \(p-1\) depends on the lift parameters
\((h_t, h_y)\). To ensure \(\nu_2(p-1) \geq s_p\), we extend the search
grid beyond the default \([-20, 20]^2\):

\[\text{half\_range} = 20 + 2^{\max(0,\, s_p - 3m)}\]

This extension guarantees sufficient coverage of bit patterns in \(p-1\)
at positions above \(3m\). For \(s_p = 36\) and \(m = 12\) (so
\(3m = 36\)), the grid remains at \([-21, 21]^2\)---only marginally
larger than the default.

\subsection{Algorithm}\label{algorithm}

The complete NTT-friendly Cocks-Pinch construction proceeds as follows:

\textbf{Input:} Embedding degree \(k=18\), target sizes \((|r|, |p|)\),
minimum two-adicity \(s\)

\textbf{Output:} Primes \((p, r)\) with \(\nu_2(r-1) \geq s\) and
\(\nu_2(p-1) \geq s\)

\begin{enumerate}
\def\labelenumi{\arabic{enumi}.}
\tightlist
\item
  Set \(m \leftarrow \lceil s/3 \rceil\) and
  \(\text{step} \leftarrow 2^m\)
\item
  Repeat:

  \begin{enumerate}
  \def\labelenumii{\alph{enumii}.}
  \tightlist
  \item
    Choose \(T\) as a random multiple of step in the target range
  \item
    Compute \(r \leftarrow T^6 - T^3 + 1\)
  \item
    If \(r\) is not prime or \(|r| \neq\) target, continue
  \item
    Compute \(\beta \leftarrow \sqrt{-3} \pmod{r}\) (Tonelli-Shanks)
  \item
    For each \(i\) with \(\gcd(i, 18) = 1\):

    \begin{itemize}
    \tightlist
    \item
      \(t_0 \leftarrow T^i + 1 \pmod{r}\)
    \item
      \(y_0 \leftarrow (t_0 - 2) \cdot \beta^{-1} \pmod{r}\)
    \item
      \(\text{range} \leftarrow 20 + 2^{\max(0, s - 3m)}\)
    \item
      For \((h_t, h_y) \in [-\text{range}, \text{range}]^2\):

      \begin{itemize}
      \tightlist
      \item
        \(t \leftarrow t_0 + h_t \cdot r\);
        \(y \leftarrow y_0 + h_y \cdot r\)
      \item
        \(p \leftarrow (t^2 + 3y^2) / 4\)
      \item
        If \(|p| =\) target AND \(\nu_2(p-1) \geq s\) AND \(p\) is
        prime: return \((p, r, t, y)\)
      \end{itemize}
    \end{itemize}
  \end{enumerate}
\item
  Until max attempts reached
\end{enumerate}

\subsection{Complexity Analysis}\label{complexity-analysis}

The following table compares the traditional and improved algorithms.
The NTT constraint increases the average number of attempts by
approximately 10--60\(\times\) (from \textasciitilde100--500 to
\textasciitilde3,000--6,000), but the absolute generation time remains
under one minute---a one-time cost amortized over the lifetime of the
curve parameters.

\begin{table}[t]
\centering
\caption{Traditional vs.~Improved Cocks-Pinch}
\begin{tabular}{lll}
\toprule
Metric & Traditional & Improved \\
\midrule
Constraint on \(T\) & None & \(T \equiv 0 \pmod{2^{12}}\) \\
\(\nu_2(r-1)\) & Random (1--3) & \(\geq 36\) (guaranteed) \\
\(\nu_2(p-1)\) & Random (1--3) & \(\geq 36\) (with ext. grid) \\
Grid size & \([-20,20]^2\) & \([-21,21]^2\) \\
Avg. attempts & 100--500 & 3,000--6,000 \\
Generation time & 2--10 s & 30--60 s \\
NTT support & No & Up to \(2^{36}\) \\
\bottomrule
\end{tabular}
\end{table}

\subsection{Comparison with BLS12-381}\label{comparison-with-bls12-381}

BLS12-381 \cite{BLS12_381}, the current industry standard for ZK proof
systems, achieves two-adicity 32 in both fields. Our construction
achieves 36, enabling NTT of length
\(2^{36} \approx 6.8 \times 10^{10}\)---a \(16\times\) improvement in
maximum polynomial degree. The tradeoff is a larger \(\rho\)-value (2.0
vs.~1.5) and correspondingly larger field elements, which is acceptable
for applications prioritizing long-term security (256-bit ECDLP
vs.~128-bit for BLS12-381).
